\scanpage{7}
\mspage{5}
car dans $S\{X_i\}[dX_i]$ l'idéal formé des formes \inserted{à p.d.} d'augmentation $\in J$ est à p.d.
\[
\big(S\{X_i\}_{0\leq i\leq n}/(\textstyle\sum X_i-t)_{\mathrm{pd}}\big)
\otimes_S\bigwedge^{\ast}\big(S^{[n]}/\mathrm{diag}\,S^{[n]}\big)
\]
C'est une $S$-algèbre diff. \inserted{graduée} anti comm., à degrés $\geq 0$, \inserted{augmentée vers $S/J$, et} à p.d. sur l'idéal \deleted{engendré par $\bigwedge^+,S\{X_i\}^+$ et $J\{X_i\}^0$ \unsure{isomorphe au}} noyau de l'augmentation
\[
C^{\bullet}_{\mathrm{DRpd}\,[n]}(S,J,t)\longrightarrow S/J\ ;
\]
% slanted
\margin{la diff est compatible avec les p.d. i.e. $d(x^{(n)})=x^{(n-1)}dx$}
et comme telle isomorphe à $S\{X_0,\ldots,X_{n-1}\}[dX_0,\ldots,dX_{n-1}\unsure{\}|\mathclose{]}}$ qui \uline{est une résolution de $S$}. \deleted{Donc a) est acquis. Mais faute de b), on \inserted{ne} peut dire en général ce \deleted{qu'il en soit} qu'est $H^{\ast}C^\ast$.}

Pour $[n]$ variable, on trouve
\[
C^{\bullet}_{\mathrm{DRpd}\,\ast}(S,J,t)=
\big(C^{\bullet}_{\mathrm{DRpd}\,[n]}(S,J,t)\big)_{n\geq 0}
\]
qui est une rés. $1/2$ simpliciale (et $=$ simpliciale) de $S$, \inserted{avec aug. vers $S/J=k$, et p.d. sur l'idéal d'augm., \unsure{compatible}} Elle dépend fonctoriellement de $(S,J,t)$, et elle permet de former, \inserted{pour $\unsure{\mathscr X|X|\mathcal X}_{\ast}\in(\mathrm{Ss})$}
\[
C^{\bullet}_{\mathrm{DRpd}}(\mathscr X_{\ast};S,J,t)=\mathrm{Hom}\big(\mathscr X_{\ast},C^{\bullet}_{\mathrm{DRpd}\,\ast}(S,J,t)\big)
\]
foncteur contravariant en $\mathscr X_{\ast}$\unsure{!} (et si $X$ esp. top
\[
\begin{aligned}
C^{\bullet}_{\mathrm{DRpd}}(X;\qquad)&=C^{\bullet}_{\mathrm{DRpd}}(S_{\ast}(X);\text{---})\\
&=\mathrm{Hom}\big(S_{\ast}(X),C^{\bullet}_{\mathrm{DRpd}\,\ast}(\quad)\big).
\end{aligned}
\]
% margin note written sideways
\margin{\unsure{!} à valeurs dans les $\unsure{S|\mathbb Z}$-algèbres graduées diff. $S/J$-augmentées à p.d. sur l'idéal d'augmentation, \unsure{compatibles|comp. tibles} avec la diff.}

Mais on ne peut dire en général \inserted{(on a seulement $H^{\bullet}_{\mathrm{DRpd}}(\mathscr X_{\ast};S,J,t)\longrightarrow H^\bullet(\mathscr X_{\ast},S)$)} quelle est sa cohomologie, et en tous cas elle ne sera pas isom (pour $\mathscr X_{\ast}$ variable) à $H^\bullet(\mathscr X_{\ast},S)$.

% archive sheet number: 6
